\documentclass[12pt, a4paper]{article}

% ── Packages ──────────────────────────────────────────────────────────────────
\usepackage[a4paper, margin=2.5cm]{geometry}
\usepackage{graphicx}
\usepackage{float}
\usepackage{amsmath}
\usepackage{amssymb}
\usepackage{xcolor}
\usepackage{mdframed}
\usepackage{parskip}
\usepackage{enumitem}
\usepackage{titlesec}
\usepackage{booktabs}
\usepackage{hyperref}


\usepackage[T1]{fontenc}

% ── Colours ───────────────────────────────────────────────────────────────────
\definecolor{keyblue}{RGB}{30, 80, 160}
\definecolor{keybluebg}{RGB}{220, 235, 255}
\definecolor{noteorange}{RGB}{200, 100, 0}
\definecolor{noteorangebg}{RGB}{255, 240, 215}
\definecolor{warnred}{RGB}{180, 30, 30}
\definecolor{warnredbg}{RGB}{255, 220, 220}
\definecolor{defgreen}{RGB}{20, 110, 50}
\definecolor{defgreenbg}{RGB}{215, 245, 225}
\definecolor{headernavy}{RGB}{25, 35, 90}

% ── Box environments ──────────────────────────────────────────────────────────

% Key Takeaway box (blue)
\newmdenv[
  linecolor=keyblue,
  backgroundcolor=keybluebg,
  linewidth=1.5pt,
  leftline=true, rightline=false, topline=false, bottomline=false,
  innerleftmargin=10pt,
  innerrightmargin=8pt,
  innertopmargin=6pt,
  innerbottommargin=6pt,
  skipabove=8pt,
  skipbelow=8pt,
]{keybox}

% Note box (orange)
\newmdenv[
  linecolor=noteorange,
  backgroundcolor=noteorangebg,
  linewidth=1.5pt,
  leftline=true, rightline=false, topline=false, bottomline=false,
  innerleftmargin=10pt,
  innerrightmargin=8pt,
  innertopmargin=6pt,
  innerbottommargin=6pt,
  skipabove=8pt,
  skipbelow=8pt,
]{notebox}

% Warning box (red)
\newmdenv[
  linecolor=warnred,
  backgroundcolor=warnredbg,
  linewidth=1.5pt,
  leftline=true, rightline=false, topline=false, bottomline=false,
  innerleftmargin=10pt,
  innerrightmargin=8pt,
  innertopmargin=6pt,
  innerbottommargin=6pt,
  skipabove=8pt,
  skipbelow=8pt,
]{warnbox}

% Definition box (green)
\newmdenv[
  linecolor=defgreen,
  backgroundcolor=defgreenbg,
  linewidth=1.5pt,
  leftline=true, rightline=false, topline=false, bottomline=false,
  innerleftmargin=10pt,
  innerrightmargin=8pt,
  innertopmargin=6pt,
  innerbottommargin=6pt,
  skipabove=8pt,
  skipbelow=8pt,
]{defbox}

% ── Section formatting ────────────────────────────────────────────────────────
\titleformat{\section}[block]
  {\large\bfseries\color{headernavy}}
  {}
  {0pt}
  {}
  [\vspace{2pt}\titlerule[0.8pt]\vspace{4pt}]

\titleformat{\subsection}[block]
  {\normalsize\bfseries\color{keyblue}}
  {}
  {0pt}
  {}

% ── Commands for labels ───────────────────────────────────────────────────────
\newcommand{\keynote}[1]{\textbf{\color{keyblue}#1}}
\newcommand{\warn}[1]{\textbf{\color{warnred}#1}}

% ── Document ──────────────────────────────────────────────────────────────────
\begin{document}

% Title
\begin{center}
  {\LARGE\bfseries\color{headernavy} Sensors \& Systems}\\[4pt]
  {\large Simple Sensors — Slides 1–9}\\[2pt]
  {\normalsize Jan Bendtsen, Torben Knudsen \quad|\quad Electronic Systems, 8th Semester}
\end{center}

\vspace{4pt}
\hrule height 1.2pt
\vspace{12pt}

% ==============================================================================
\section*{Slides 1–3 \quad What Is a ``Simple'' Sensor?}

A \textbf{simple sensor} is a self-contained device that converts a single,
continuous physical quantity into an electrical signal that can be read
directly with little or no additional computation.
Three defining properties hold for all simple sensors:

\begin{itemize}[leftmargin=1.5em]
  \item It measures \textbf{one continuous physical quantity} — temperature,
        magnetic field strength, sound pressure, humidity, etc.
  \item It provides a \textbf{simple interface} — either an analog voltage/current
        output or a straightforward digital protocol (I\textsuperscript{2}C,
        SPI, UART, etc.).
  \item It does \textbf{not require extensive processing} to obtain the desired
        measurement — contrast this with a camera or radar, where significant
        computation must precede any useful output.
\end{itemize}

Examples encountered in this lecture include an air-quality sensor, a Hall
effect sensor, and a microphone module — all pictured on slide~3.
Despite measuring entirely different physical quantities, every one of
them fits the same three-property mold.

\begin{figure}[H]
  \centering
  \includegraphics[width=0.70\linewidth]{figures/slide-03.jpg}
  \caption{Three ``simple'' sensors: air-quality (left), Hall effect (centre),
           microphone (right). Source: lecture slides.}
  \label{fig:simple_sensors}
\end{figure}

\begin{defbox}
  \textbf{Definition — Simple Sensor:} a sensor that
  (1)~measures a single continuous physical quantity,
  (2)~provides a simple analog or digital interface, and
  (3)~requires no extensive post-processing to obtain the desired
  information.
\end{defbox}

% ==============================================================================
\section*{Slide 4 \quad The Sensor Signal Chain}

Every sensor — simple or otherwise — follows the same general processing
pipeline from the physical world to a number in your microcontroller.
Understanding this chain is important because it tells you \emph{where} a
particular sensor sits and what processing remains your responsibility.

\begin{figure}[H]
  \centering
  \includegraphics[width=0.95\linewidth]{figures/slide-04.jpg}
  \caption{The sensor signal chain: physical event $\to$ sensing component
           $\to$ amplifier $\to$ ADC $\to$ computing.}
  \label{fig:signal_chain}
\end{figure}

Walking through each stage:

\begin{enumerate}[leftmargin=1.5em]
  \item \textbf{Physical event.} A real-world phenomenon occurs — light,
        sound, temperature change, force, radiation. This is the stimulus
        the sensor exists to measure.

  \item \textbf{Sensing component.} A transducer converts the physical
        phenomenon into an electrical signal, typically a voltage or current.
        This is the core physics step.  The output is often tiny — millivolts
        or microamps.

  \item \textbf{Amplifier.} The weak electrical signal is boosted to a level
        suitable for the ADC. Because sensor signals are small, amplifier
        quality matters; any noise or offset introduced here gets amplified
        alongside the signal.

  \item \textbf{Analog-to-Digital Converter (ADC).} The amplified continuous
        signal is sampled at regular intervals $T_s$ and converted to a
        digital integer. This is where the continuous physical world becomes a
        discrete-time sequence of values. Output protocols include
        PWM, RS232, I\textsuperscript{2}C, SPI, and CAN.

  \item \textbf{Computing.} A microcontroller or processor interprets the
        stream of digital values and acts on them.
\end{enumerate}

\begin{notebox}
  \textbf{Note — integration level.} Different sensors integrate different
  numbers of these stages internally.  A bare PT100 resistance thermometer
  gives you only a raw resistance — stages 2 onward are your problem.
  A TMP117 integrates the amplifier and ADC on-chip and hands you
  temperature in \textsuperscript{o}C directly over I\textsuperscript{2}C.
  Always check the datasheet to know what the sensor does for you.
\end{notebox}

\begin{keybox}
  \textbf{Key Takeaway:} The same five-stage chain governs every sensor
  system.  ``Simple'' sensors merely make stages 3–4 (and sometimes 5)
  invisible to the user by integrating them on-chip.
\end{keybox}

% ==============================================================================
\section*{Slides 5–7 \quad The Hall Effect Sensor}

\subsection*{What Is a Hall Sensor For?}

A Hall sensor is used to \textbf{measure or detect magnetic fields}.
Just as a thermistor converts temperature into a resistance that you can
read, a Hall sensor converts magnetic field strength (flux density,
measured in Gauss or Tesla) into a voltage that you can read.

Practical uses include:
\begin{itemize}[leftmargin=1.5em]
  \item Detecting the presence of a magnet (e.g.\ reed-switch replacement,
        door/lid open-closed detection).
  \item Measuring rotation speed — a magnet on a shaft passes the sensor
        once per revolution, producing a pulse.
  \item Measuring current — a current-carrying conductor generates a
        magnetic field proportional to the current.
  \item Measuring the position of a magnetic actuator or joystick.
\end{itemize}

\subsection*{The Physics of the Hall Effect (Slide 6)}

\begin{figure}[H]
  \centering
  \includegraphics[width=0.95\linewidth]{figures/slide-06.jpg}
  \caption{Left: no magnetic field, $V_\mathrm{HALL}\approx 0$\,V.
           Right: magnetic field perpendicular to bias current,
           $V_\mathrm{HALL}$ rises proportionally.}
  \label{fig:hall_physics}
\end{figure}

Inside the sensor there is a thin \textbf{semiconductor sheet}.
Tracing through the signal chain of Figure~\ref{fig:signal_chain} for
this specific sensor:

\textbf{Stage 1 — Physical event:} A magnetic field is present (or not)
in the vicinity of the sensor.

\textbf{Stage 2 — Sensing component (the Hall element):} A constant
voltage source forces a fixed \textbf{bias current} $I_\mathrm{BIAS}$
through the semiconductor sheet along its length (the element marked
with~$\times$ in the left diagram).

\begin{itemize}[leftmargin=1.5em]
  \item \textbf{No field (left diagram):} charge carriers flow straight
        through the sheet; the voltage measured \emph{across the width}
        of the sheet, $V_\mathrm{HALL}$, is approximately zero.
  \item \textbf{Field present (right diagram):} flux lines enter the
        semiconductor perpendicular to $I_\mathrm{BIAS}$ (the green
        arrows passing through the block).  The Lorentz force
        $\mathbf{F} = q\mathbf{v} \times \mathbf{B}$ deflects the moving
        charge carriers sideways, causing them to accumulate along one
        edge of the sheet.  This charge separation creates a measurable
        transverse voltage — $V_\mathrm{HALL}$ rises toward a positive
        value.
\end{itemize}

The key relationship is:
\[
  V_\mathrm{HALL} \propto I_\mathrm{BIAS} \cdot B
\]
where $B$ is the magnetic flux density.  Since $I_\mathrm{BIAS}$ is
kept constant by the circuit, $V_\mathrm{HALL}$ varies \textbf{linearly}
with $B$ — stronger field, larger voltage.

\begin{defbox}
  \textbf{Variable summary for the Hall element:}
  \begin{itemize}[leftmargin=1.5em, topsep=2pt]
    \item $I_\mathrm{BIAS}$ — constant bias current forced through the
          semiconductor sheet; set by the external voltage source.
    \item $B$ — magnetic flux density of the external field
          (Gauss\,/\,Tesla); the quantity being measured.
    \item $V_\mathrm{HALL}$ — output voltage measured across the
          \emph{width} of the sheet; approximately zero with no field,
          proportional to $B$ when a field is present.
  \end{itemize}
\end{defbox}

\subsection*{Signal Conditioning — Why Each Configuration Differs (Slide 7)}

\begin{figure}[H]
  \centering
  \includegraphics[width=0.95\linewidth]{figures/slide-07.jpg}
  \caption{Left: linear Hall sensor (Hall element + OpAmp only).
           Right: switch-type Hall sensor (Hall element + OpAmp +
           comparator/Schmitt trigger).}
  \label{fig:hall_conditioning}
\end{figure}

$V_\mathrm{HALL}$ is typically only a few millivolts — far too small to
be useful on its own.  This brings us to \textbf{Stage 3 (Amplifier)}
in the signal chain.  Two different goals lead to two different circuit
configurations:

\subsubsection*{Configuration 1 — Measuring Field Strength (left circuit)}

If you need to know \emph{how strong} the magnetic field is, you connect
a precision operational amplifier (OpAmp) directly after the Hall element.
The OpAmp amplifies $V_\mathrm{HALL}$ to a usable voltage level, and the
analog output is proportional to $B$.  In terms of the signal chain, you
stop after Stage~3 (or pass to an ADC for Stage~4); you keep the output
continuous and linear.

\begin{warnbox}
  \warn{Warning:} this configuration demands a \textbf{high-quality,
  low-offset OpAmp}.  Any DC offset in the OpAmp is amplified just like
  the signal — a 1\,mV offset with a gain of 500 becomes 500\,mV of
  error, completely swamping a small magnetic signal.
\end{warnbox}

\subsubsection*{Configuration 2 — Detecting Presence Only (right circuit)}

If you only need to know \emph{whether} a magnet is nearby — a binary
yes/no — you add a \textbf{comparator (Schmitt trigger)} after the
amplifier.  The comparator compares the amplified $V_\mathrm{HALL}$
against a threshold and outputs a clean digital HIGH or LOW.

Why is this circuit easier and cheaper?
\begin{itemize}[leftmargin=1.5em]
  \item The comparator threshold absorbs small OpAmp offsets — the
        binary output is correct as long as the field is clearly above
        or below the threshold.
  \item A lower-grade OpAmp is acceptable because small gain errors do
        not affect the final digital output.
  \item The Schmitt trigger adds \textbf{hysteresis}: the turn-on and
        turn-off thresholds differ slightly, preventing rapid chattering
        when the magnet is near the boundary.
\end{itemize}

\begin{keybox}
  \textbf{Key Takeaway — choosing a configuration:}
  \begin{itemize}[leftmargin=1.5em, topsep=2pt]
    \item Need a \emph{proportional} measurement of field strength?
          $\to$ Linear Hall sensor (OpAmp only).
          Requires a high-quality OpAmp; output is analog.
    \item Need only a \emph{presence/absence} detection?
          $\to$ Switch-type Hall sensor (OpAmp + comparator).
          Tolerates lower-grade OpAmp; output is digital.
  \end{itemize}
  The two configurations differ because they have different goals — the
  stopping point in the signal chain (Stage~3 vs.\ a binary Stage~4)
  drives everything else.
\end{keybox}

% ==============================================================================
\section*{Slide 8 \quad The SS49E Linear Hall Sensor}

\begin{figure}[H]
  \centering
  \includegraphics[width=0.55\linewidth]{figures/slide-08.jpg}
  \caption{The SS49E linear Hall effect sensor module.}
  \label{fig:ss49e}
\end{figure}

The SS49E is a specific \textbf{linear} Hall sensor — it uses
Configuration~1 above and provides a continuous analog output
proportional to field strength.  Key specifications:

\begin{itemize}[leftmargin=1.5em]
  \item \textbf{Sensitivity: $2.5 \pm 0.5$\,mV/G.}
        For every Gauss of magnetic flux density, the output voltage
        changes by approximately 2.5\,mV.  The $\pm 0.5$\,mV/G
        tolerance means unit-to-unit sensitivity variation of up to 20\%
        — calibration is advisable in precision applications.
  \item \textbf{Power consumption: 4\,mA at 5\,V DC.}
        Modest consumption; compatible with direct Arduino/microcontroller
        supply.
  \item \textbf{Minimum supply: 2.3\,V.}
        The sensor can run from 3.3\,V systems (Raspberry Pi, Nordic MCUs,
        etc.) without a level shifter.
\end{itemize}

\begin{notebox}
  \textbf{Note — sensitivity tolerance.} The $\pm 0.5$\,mV/G spread
  means two sensors from the same batch may give readings that differ
  by up to 20\% for the same field.  For precise measurements, each
  unit should be individually calibrated against a known field.
\end{notebox}

% ==============================================================================
\section*{Slide 9 \quad Arduino Example — Reading the SS49E}

\begin{figure}[H]
  \centering
  \includegraphics[width=0.95\linewidth]{figures/slide-09.jpg}
  \caption{Wiring diagram and Arduino code for reading the SS49E.}
  \label{fig:hall_code}
\end{figure}

The example code on slide~9 demonstrates a complete reading pipeline.
Three steps worth unpacking:

\subsubsection*{Step 1 — Averaging to reduce noise}
\begin{verbatim}
long measure = 0;
for(int i = 0; i < 10; i++){
    measure += analogRead(pinHall);
}
measure /= 10;
\end{verbatim}

Ten successive ADC readings are accumulated and then divided by 10.
This is a basic \textbf{averaging filter} — a first encounter with a
theme that runs throughout this lecture.  Averaging $N$ independent
noise samples reduces the noise standard deviation by a factor of
$\sqrt{N}$, so averaging 10 readings gives $\sqrt{10}\approx 3.2\times$
noise reduction at the cost of 10 ADC conversions.

\subsubsection*{Step 2 — ADC count to millivolts}
\begin{verbatim}
float outputV = measure * 5000.0 / 1023;
\end{verbatim}

The Arduino Uno's ADC has \textbf{10-bit resolution}: it maps
0–5\,V onto the integer range 0–1023.  Multiplying by $5000/1023$
undoes this mapping and recovers the voltage in millivolts.

\subsubsection*{Step 3 — Millivolts to magnetic flux density}
\begin{verbatim}
float magneticFlux = outputV * 53.33 - 133.3;
\end{verbatim}

This is the \textbf{inverse of the sensor's sensitivity curve},
converting millivolts back to flux density in millitesla.  The
coefficients 53.33 and 133.3 come from the SS49E datasheet.
The subtraction of 133.3 accounts for the sensor's quiescent output
voltage at zero field (the output sits at mid-supply, not at 0\,V, when
no field is present — so zero field must be subtracted out).

\begin{keybox}
  \textbf{Key Takeaways — Hall sensor end-to-end:}
  \begin{itemize}[leftmargin=1.5em, topsep=2pt]
    \item A Hall sensor converts \textbf{magnetic field strength} to a
          voltage via the Hall effect.
    \item $V_\mathrm{HALL}$ is tiny; an OpAmp is always needed.
    \item For field \emph{strength} measurements use a linear sensor
          (analog output). For field \emph{presence} detection use a
          switch-type sensor (digital output).
    \item The SS49E has a sensitivity of $2.5 \pm 0.5$\,mV/G;
          calibration is recommended for precision work.
    \item Averaging multiple ADC readings is a simple and effective
          way to reduce measurement noise — a preview of the filtering
          concepts covered later in this lecture.
  \end{itemize}
\end{keybox}

\end{document}
